%% =========================================================
%% Section: the categorical anomaly of the Fibonacci symmetry as a
%% rank-one no-go, continuing the fusion examples of
%% chapter/ch25_asymmetric_examples.tex, Section~\ref{sec:asymex_fibonacci}.
%% =========================================================

\section{The Fibonacci categorical anomaly}%
\label{sec:asymex_fib_anomaly}

The Fibonacci fusion ring $\Z[\tau]/(\tau^2-\tau-1)$ has no unital ring
homomorphism to $\Z$: the two possible images of $\tau$ over $\C$,
$\varphi=(1+\sqrt5)/2$ and $-1/\varphi$, are both irrational.
In the theory of matrix product operator symmetries of
Ref.~\cite{GarreRubio2023MPOSym}, the multiplicities with which the normal
blocks of a symmetric matrix product state reappear under the symmetry
are nonnegative integers forming a representation of the fusion
rules~\cite[line~683]{GarreRubio2023MPOSym}. The Fibonacci example has two
blocks $x_1,x_\tau$ with $\tau\cdot x_1=x_\tau$ and
$\tau\cdot x_\tau=x_1+x_\tau$~\cite[lines~1991--1993]{GarreRubio2023MPOSym}.
This section realizes that regular representation and proves that no
single normal tensor of positive bond dimension can carry the symmetry
with length-independent scalar action and with the unit acting as the
identity. The scalar would have to be a nonnegative integer satisfying
$m^2=m+1$. This obstruction does not classify representations with more
than one block; in particular, it does not exclude direct sums of the
regular representation. It is a non-invertible analogue of the obstruction
to an invariant single-block state from a non-trivial
three-cocycle~\cite[line~738]{GarreRubio2023MPOSym}, not a computation of a
three-cocycle associator.

\begin{definition}[The Fibonacci fusion matrices]%
    \label{def:asymex_fib_anomaly_nim}
    \lean{FibonacciCompression.fibFusionMatrix, FibonacciCompression.fibNim,
        FibonacciCompression.fibBlock}
    \leanok
    \uses{def:asymex_fib}
    Index the two simple objects of Definition~\ref{def:asymex_fib} by
    $0$ and $1$, corresponding respectively to the trivial object $1$
    and the object $\tau$. The fusion matrix
    \begin{align}
        N_\tau & =\begin{pmatrix}0&1\\1&1\end{pmatrix}
        \label{eq:asymex_fib_anomaly_matrix}
    \end{align}
    represents multiplication by $\tau$ in the basis $1,\tau$ of the
    fusion ring. Together with $N_1=\Id_2$, it gives the regular
    representation by matrices with nonnegative integer entries.
    For $a,b,c\in\{1,\tau\}$, let $N_{ab}^c$ be the $(b,c)$ entry of
    $N_a$. The corresponding matrix product operator tensors are the
    admissibility block $B_1$ and the block $B_\tau$ of
    Definition~\ref{def:asymex_fib}.
\end{definition}

\begin{theorem}[The Fibonacci relation for the fusion matrix]%
    \label{thm:asymex_fib_anomaly_matrix_sq}
    \lean{FibonacciCompression.fibFusionMatrix_sq}
    \leanok
    \uses{def:asymex_fib_anomaly_nim}
    $N_\tau^2=N_\tau+\Id_2$, as matrices with nonnegative integer entries.
\end{theorem}

\begin{proof}
    \leanok
    Direct multiplication gives
    $N_\tau^2=\begin{pmatrix}1&1\\1&2\end{pmatrix}=N_\tau+\Id_2$.
\end{proof}

\begin{theorem}[The fusion ring on the periodic operators]%
    \label{thm:asymex_fib_anomaly_fusion_algebra}
    \lean{FibonacciCompression.fibonacci_fusion_algebra}
    \leanok
    \uses{def:asymex_fib_anomaly_nim, thm:asymex_fib_fusion}
    For every pair of labels $a,b\in\{1,\tau\}$ and every positive length
    $L$, the periodic operators $O_a=O_a(L)=O_L(B_a)$ satisfy
    \begin{align}
        O_aO_b & =\sum_cN_{ab}^cO_c,
        \label{eq:asymex_fib_anomaly_fusion_algebra}
    \end{align}
    that is, $O_1O_1=O_1$, $O_1O_\tau=O_\tau O_1=O_\tau$, and
    $O_\tau O_\tau=O_1+O_\tau$.
    \begin{center}
        % Formula/source: the (tau,tau) and (1,tau) cases of
        % (\ref{eq:asymex_fib_anomaly_fusion_algebra}), i.e.
        % O_tau O_tau = O_1 + O_tau (already (\ref{eq:asymex_fib_fusion}))
        % and O_1 O_tau = O_tau.
        % Ink-to-index map: each row of L tensors is the periodic operator
        % O_a(L), built from L copies of the block B_a; the top row's
        % physical legs i_1,...,i_L are the free outgoing labels, the
        % bottom row's physical legs j_1,...,j_L are the free incoming
        % labels, and the unlabelled physical legs shared between the two
        % rows are the L contracted intermediate labels of the
        % composition O_a(L)*O_b(L).
        % Contraction set: the L intermediate physical legs between the
        % two rows, and, independently on each row, the periodic virtual
        % bond of that row's block.
        % Boundary signature: (virtual-west=0 | virtual-east=0 |
        % physical-up=L | physical-down=L); no virtual legs remain open,
        % both rows being closed by their own periodic trace.
        \tenkzkernel
        \begin{tenkz}[rows={op,op}, cols=4, boundary=periodic]
            \tn[skin=mpo, ports={90:physical:$i_1$, 270:physical}]{B_\tau} &
            \tn[skin=mpo, ports={90:physical:$i_2$, 270:physical}]{B_\tau} &
            \tn[skin=dots, ports={180:virtual, 0:virtual}]{} &
            \tn[skin=mpo, ports={90:physical:$i_L$, 270:physical}]{B_\tau} \\
            \tn[skin=mpo, ports={90:physical, 270:physical:$j_1$}]{B_\tau} &
            \tn[skin=mpo, ports={90:physical, 270:physical:$j_2$}]{B_\tau} &
            \tn[skin=dots, ports={180:virtual, 0:virtual}]{} &
            \tn[skin=mpo, ports={90:physical, 270:physical:$j_L$}]{B_\tau}
        \end{tenkz}
        \qquad\text{and}\qquad
        \begin{tenkz}[rows={op,op}, cols=4, boundary=periodic]
            \tn[skin=mpo, ports={90:physical:$i_1$, 270:physical}]{B_1} &
            \tn[skin=mpo, ports={90:physical:$i_2$, 270:physical}]{B_1} &
            \tn[skin=dots, ports={180:virtual, 0:virtual}]{} &
            \tn[skin=mpo, ports={90:physical:$i_L$, 270:physical}]{B_1} \\
            \tn[skin=mpo, ports={90:physical, 270:physical:$j_1$}]{B_\tau} &
            \tn[skin=mpo, ports={90:physical, 270:physical:$j_2$}]{B_\tau} &
            \tn[skin=dots, ports={180:virtual, 0:virtual}]{} &
            \tn[skin=mpo, ports={90:physical, 270:physical:$j_L$}]{B_\tau}
        \end{tenkz}
    \end{center}
    The first diagram, closed as a periodic chain of length $L$, is
    $O_\tau(L)O_\tau(L)$; the second is $O_1(L)O_\tau(L)$.
\end{theorem}

\begin{proof}
    \leanok
    The case $a=b=1$ is the admissibility projector composed with itself.
    The cases $\{a,b\}=\{1,\tau\}$ are an explicit block-triangular
    compression computation over $\Z[\sigma]$, transported to the complex
    matrices as in Theorem~\ref{thm:asymex_golden_ring}, in the style of
    Theorem~\ref{thm:asymex_fib_compression}. The case $a=b=\tau$ is
    Theorem~\ref{thm:asymex_fib_fusion}.
\end{proof}

\begin{theorem}[The regular representation on a pair of normal states]%
    \label{thm:asymex_fib_anomaly_nim_rep}
    \lean{FibonacciCompression.fibonacci_nim_rep,
        FibonacciCompression.mpo_fibOne_mulVec_chain}
    \leanok
    \uses{def:asymex_fib_anomaly_nim, thm:asymex_fib_anomaly_fusion_algebra,
        cor:asym_action_tensors}
    There is a pair of normal matrix product state tensors of positive
    bond dimension, with periodic vectors $\Psi_1,\Psi_\tau$, realizing the
    all-$\tau$ product state and its image under $O_\tau$
    (Corollary~\ref{cor:asym_action_tensors}), such that at every positive
    length $L$,
    \begin{align}
        O_a\Psi_s & =\sum_tN_{as}^t\Psi_t,
        \label{eq:asymex_fib_anomaly_nim_rep}
    \end{align}
    that is, $O_\tau\Psi_1=\Psi_\tau$, $O_\tau\Psi_\tau=\Psi_1+\Psi_\tau$,
    and $O_1$ fixes both $\Psi_1$ and $\Psi_\tau$. Here the dependence of
    the periodic vectors on $L$ is suppressed. The multiplicities are
    the rows of the fusion matrix (\ref{eq:asymex_fib_anomaly_matrix}).
    \begin{center}
        % Formula/source: the sitewise fusion tensors underlying
        % (\ref{eq:asymex_fib_anomaly_nim_rep}), specializing
        % Theorem~thm:asymex_fib_fusion_tensors to the slots s=1,tau; the
        % remainder of the compression of Theorem~thm:asymex_fib_compression
        % vanishes identically, so the intertwining identities below hold
        % exactly, with no residual term.
        % Ink-to-index map: B is the stacked source tensor of
        % Definition~def:asymex_fib, with physical leg i; V_s is the
        % column block of the adapted basis carrying the slot s, and W_s
        % the matching row block of its inverse; C_s is the block of the
        % target of label s.
        % Contraction set: only the wire joining B (or C_s) to V_s (or
        % W_s) is summed; no other index is contracted.
        % Boundary signature: (virtual-west=1 | virtual-east=1 |
        % physical-up=1 | physical-down=0) on every wire.
        \tenkzkernel
        \begin{tenkz}[rows={wire}, cols=2, west=open, east=open]
            \tn[label pos=s, ports={90:physical:$i$}]{B} & \tn[skin=ring]{V_s}
        \end{tenkz}
        \quad $=$ \quad
        \begin{tenkz}[rows={wire}, cols=2, west=open, east=open]
            \tn[skin=ring]{V_s} & \tn[label pos=s, ports={90:physical:$i$}]{C_s}
        \end{tenkz}
        \\[2ex]
        \begin{tenkz}[rows={wire}, cols=2, west=open, east=open]
            \tn[skin=ring]{W_s} & \tn[label pos=s, ports={90:physical:$i$}]{B}
        \end{tenkz}
        \quad $=$ \quad
        \begin{tenkz}[rows={wire}, cols=2, west=open, east=open]
            \tn[label pos=s, ports={90:physical:$i$}]{C_s} & \tn[skin=ring]{W_s}
        \end{tenkz}
    \end{center}
\end{theorem}

\begin{proof}
    \leanok
    \uses{thm:asymex_fib_fusion_tensors}
    The state $\Psi_1$ is the all-$\tau$ product state and $\Psi_\tau$ is
    its image under the periodic operator of the $\tau$ block. The case
    $a=\tau$ is the fusion-tensor identity of
    Theorem~\ref{thm:asymex_fib_fusion_tensors} applied to the source
    stacked square. The cases $a=1$ reduce, by an explicit compression
    computation over $\Z[\sigma]$ analogous to
    Theorem~\ref{thm:asymex_fib_anomaly_fusion_algebra}, to the unit law
    fixing each of $\Psi_1,\Psi_\tau$.
\end{proof}

\begin{theorem}[Word-trace relations against a normal tensor are integral]%
    \label{thm:asymex_fib_anomaly_trace_integral}
    \lean{MPSTensor.exists_nat_eq_of_forall_trace_evalWord_eq_mul}
    \leanok
    \uses{def:normal}
    Let $B$ be a matrix product state tensor and $A$ a normal tensor of
    positive bond dimension, over the same alphabet. If a scalar $c$
    satisfies $\tr(B^w)=c\,\tr(A^w)$ for every nonempty word $w$, then $c$
    is a nonnegative integer.
\end{theorem}

\begin{proof}
    \leanok
    Normality writes a diagonal matrix unit of $A$ as a combination
    $\sum_ux_uA^u$ of words of one positive length $\ell$. The same
    combination $Q=\sum_ux_uB^u$ then satisfies
    $\tr(Q^k)=c\,\tr((\sum_ux_uA^u)^k)=c$ for every positive $k$.
    A constant trace of all positive powers of a fixed matrix is a
    nonnegative integer.
\end{proof}

\begin{theorem}[Rank-one actions on a normal tensor are integral]%
    \label{thm:asymex_fib_anomaly_integral_action}
    \lean{MPOTensor.exists_nat_eq_of_mpo_mulVec_mpv_eq_smul}
    \leanok
    \uses{thm:asymex_fib_anomaly_trace_integral, cor:asym_action_tensors}
    Let $T$ be a matrix product operator tensor and $A$ a normal matrix
    product state tensor of positive bond dimension. If at every positive
    length the periodic operator of $T$ acts on the periodic vector of $A$
    as multiplication by one length-independent scalar $c$, then $c$ is a
    nonnegative integer.
\end{theorem}

\begin{proof}
    \leanok
    The periodic-vector relation is the word-trace relation of
    Theorem~\ref{thm:asymex_fib_anomaly_trace_integral} for the action
    tensor $T\cdot A$ of Corollary~\ref{cor:asym_action_tensors} against
    $A$.
\end{proof}

\begin{theorem}[A nonzero length-independent eigenvalue is impossible]%
    \label{thm:asymex_fib_anomaly_eigenvalue_zero}
    \lean{FibonacciCompression.fibTau_eigenvalue_eq_zero}
    \leanok
    \uses{thm:asymex_fib_anomaly_fusion_algebra,
        thm:asymex_fib_anomaly_integral_action}
    Let $A$ be a normal tensor of positive bond dimension, with periodic
    vectors $\Psi_A=\ket{V^{(L)}(A)}$. If the periodic operator of the
    $\tau$ block acts on $\Psi_A$ at every positive length by one
    length-independent scalar $c$, $O_\tau\Psi_A=c\Psi_A$, then $c=0$.
\end{theorem}

\begin{proof}
    \leanok
    By Theorem~\ref{thm:asymex_fib_anomaly_integral_action}, applied to the
    action of the $\tau$ block on $A$, $c$ is a nonnegative integer $m$. If
    $m\neq0$,
    cancelling $m$ in the relation obtained by applying the unit law
    $O_1O_\tau=O_\tau$ (Theorem~\ref{thm:asymex_fib_anomaly_fusion_algebra})
    to $\Psi_A$ gives $O_1\Psi_A=\Psi_A$. The fusion rule
    $O_\tau O_\tau=O_1+O_\tau$ applied to $\Psi_A$ then gives
    $m^2\Psi_A=(m+1)\Psi_A$. At length $\ell+1$, where $\ell$ is an
    injectivity length of $A$, injectivity ensures that the periodic
    vector is nonzero. Hence $m^2=m+1$, which no nonnegative integer
    satisfies. Therefore $m=0$.
\end{proof}

\begin{theorem}[No normal matrix product state is Fibonacci symmetric]%
    \label{thm:asymex_fib_anomaly_no_go}
    \lean{FibonacciCompression.not_exists_normal_fibonacci_symmetric}
    \leanok
    \uses{thm:asymex_fib_anomaly_eigenvalue_zero, thm:asymex_fib_fusion}
    There is no positive bond dimension $D$, normal tensor $A$ of bond
    dimension $D$, and scalar $c$, such that at every positive length $L$
    the periodic operators satisfy $O_1\Psi_A=\Psi_A$ and
    $O_\tau\Psi_A=c\Psi_A$. Thus no normal matrix product state of
    positive bond dimension realizes a one-dimensional representation of
    the Fibonacci fusion algebra with length-independent scalar action
    and with $O_1$ acting as the identity. The obstruction is the absence
    of a nonnegative integer $m$ satisfying $m^2=m+1$.
    The pair of normal states of
    Theorem~\ref{thm:asymex_fib_anomaly_nim_rep} shows that two blocks
    suffice, so the minimal number of blocks of a Fibonacci-symmetric
    family of normal states with length-independent multiplicities is two.
\end{theorem}

\begin{proof}
    \leanok
    If such $A$, $c$ existed, Theorem~\ref{thm:asymex_fib_anomaly_eigenvalue_zero}
    would give $c=0$, so $O_\tau\Psi_A=0$. The fusion rule
    $O_\tau O_\tau=O_1+O_\tau$ applied to $\Psi_A$ would then give
    $0=O_1\Psi_A=\Psi_A$ at every positive length. This contradicts the
    nonvanishing of the periodic vectors of a normal tensor of positive
    bond dimension at sufficiently large lengths.
\end{proof}
